\documentclass[../main.tex]{subfiles}
\begin{document}
%==========================================TIÊU ĐỀ TẬP========================================
\begin{table}[h]
  \begin{adjustwidth}{-1cm}{}
    \begin{tabular}{>{\centering\arraybackslash}p{6cm}|>{\raggedright\arraybackslash}p{12cm}}
      \multirow{5}{6cm}
      % Cột bên trái: ÉPISODE và số tập
      {\\[-40pt]
        \centering
        {\fontsize{20pt}{20pt}\selectfont{\outlinetext[1]{eptype}{white}{\flaregothic ÉPISODE}}\\[12pt]
          {\fontsize{72pt}{72pt}\selectfont{\outlinetext[1]{epnum}{white}{\burbank 170}}}\\[6pt]       % use for normal/special episodes
          % {\fontsize{36pt}{36pt}\selectfont{\outlinetext[1]{epnum}{white}{\burbank SÁU}}}\\[6pt]     % use for custom episodes
          {\fontsize{20pt}{20pt}\selectfont{\outlinetext[1]{epnum}{white}{\cafeteria (13-11)}}}\color{black}}
        \\[18pt]
        % {\fontsize{14pt}{14pt}\selectfont{\outlinetext[1]{epattr}{white}{\flaregothic [I]}}}\\
        % {\fontsize{18pt}{18pt}\selectfont{\outlinetext[1]{epattr}{white}{\flaregothic [II]}}}\\
        \color{black}
      }
       &
      % Dòng thứ nhất: tên tập tiếng Nhật
      {\fotskip\fontsize{18pt}{18pt}\selectfont \ruby{台}{たい}\ruby{風}{ふう}\ruby{襲}{しゅう}\ruby{来}{らい}(2022\ruby{個}{こ})で\ruby{旅}{りょ}\ruby{行}{こう}が\ruby{台}{だい}\ruby{無}{な}し} \\[6pt]
      % Dòng thứ hai: tên tập tiếng Anh
       & {\cafeteria\fontsize{18pt}{18pt}\selectfont Typhoons Ruin Trip}                                                                                              \\[4pt]
      % Dòng thứ ba: tên tập tiếng Việt
       & {\vnmsans\fontsize{18pt}{18pt}\selectfont Chuyến đi bão táp}                                                                                                 \\[8pt]
      % Dòng thứ tư: Nhân vật xuất hiện
       & {\caviar{\fontsize{14pt}{14pt}\selectfont Nhân vật: \emp{kuon2} \emp{naomi2} \emp{game} \emp{subaru} \emp{suisei} \emp{luna}}}                               \\
      % Dòng thứ năm (nếu có): Nhân vật chỉ xuất hiện trong Omake
      %  & {\abv Truyện phụ:} \emp{okayu}
      \\[8pt]
      % Dòng cuối cùng: Thời gian diễn ra của tập (thường sẽ là ngày phát hành)
       & {\sen\fontsize{16pt}{16pt}\selectfont Mốc thời gian: 14/08/2022}                                                                                             \\[18pt]
    \end{tabular}\\[8pt]
  \end{adjustwidth}
\end{table}
%=========================================TRUYỆN CHÍNH========================================
\color{black}

\txtbx[2]{
  {\color{vol} \uline{Tóm tắt:}} Chuyến du lịch biển cuối tuần của Kuon, cô em út Naomi cùng bộ ba Subaru, Suisei và Luna bỗng chốc biến thành thảm họa thiên tai khi họ đụng độ siêu bão năm 2022 đang sản sinh ra chuỗi liên hoàn 2022 cơn lốc xoáy càn quét khắp các hòn đảo. Bất chấp hiểm nguy, Luna vẫn kiên quyết sống ảo mọi lúc mọi nơi để rồi liên tục lăn ra ăn vạ khi điện thoại mất sóng. Trong khi Kuon chật vật che chắn bảo vệ em gái Naomi và cô bé liên tục bấm máy ghi lại những khoảnh khắc bi hài, cả nhóm liên tục bị các cơn lốc xoáy cuốn phăng từ đảo này sang đảo khác trước khi đáp cánh an toàn về lại văn phòng theo một kịch bản không ai ngờ tới.
}

\pov{Hikawa Kuon}

Một ngày trung tuần tháng Tám năm 2022.

Theo bản tin dự báo thời tiết khẩn cấp mà chúng tôi\dots\ quên béng không thèm kiểm tra trước khi khởi hành, một siêu bão nhiệt đới quy mô bất thường đang trên đường đổ bộ trực diện vào đất liền Thành phố Điện tử.

Và tại một hòn đảo du lịch ngoài khơi xa, nơi mà tôi, cô em gái út Naomi cùng bộ ba Suisei, Subaru và Luna vừa mới hào hứng đặt chân tới, tất cả chúng tôi đều phải chết lặng trong nỗi bàng hoàng tột độ.

Trước mắt chúng tôi là một khung cảnh tận thế thực sự: biển động gầm rú dữ dội, những đợt sóng bạc đầu cao ngút ngàn đập ầm ầm vào bờ đá, mây đen cuồn cuộn vần vũ che kín cả vòm trời. Mưa trút xuống xối xả trắng xóa cả tầm mắt, mọi cảnh quan tươi đẹp thường ngày bỗng chốc tối sầm lại. Chẳng còn lấy một chút màu xanh biếc của biển cả, cũng chẳng còn bờ cát vàng óng ả dưới ánh nắng hè rực rỡ như trong những cuốn cẩm nang du lịch quảng cáo.

Tôi vội vã kích hoạt chiếc khiên điện tử phòng hộ phát ra từ đồng hồ thông minh để che mưa cho cả nhóm. Cô em út Naomi đứng nép sát vào người tôi, hai tay ôm khư khư chiếc máy ảnh du lịch nhỏ đeo trước ngực, đôi mắt tròn xoe ngơ ngác nhìn cảnh tượng thiên nhiên thịnh nộ.

Trong khi đó, ở ngay sát mép nước, nàng Công chúa Kẹo Himemori Luna đã quỳ sụp cả hai gối xuống bờ cát ướt sũng. Cô nàng giàn giụa nước mắt, liên tục lấy hai tay đập bồm bộp xuống cát, khóc lóc ăn vạ một cách đầy đau khổ trong khi Suisei và Subaru đứng trố mắt, há hốc mồm kinh ngạc đến mức không thốt nên lời.

\img{13-11a}

Tôi nhìn cô nàng tóc hồng đang vật vã dưới mưa mà không khỏi day day trán, buông một tiếng thở dài thườn thượt: ``Tháng trước đi nghỉ mát ở bãi biển cùng gia đình anh và Gen 4 vừa mới nếm mùi bão biển kinh hoàng, cây dừa gãy đôi suýt nữa cuốn bay chiếc phao vịt của em chưa đủ tởn hay sao mà mới tròn một tháng sau, em lại rủ rê mọi người đi biển đúng ngày có siêu bão thế này hả Luna\dots?''

Luna vừa gạt nước mắt vừa ngước lên phụng phịu, giọng ngọng nghịu ăn vạ: ``Nhưng mừ hồi ở bờ biển tháng trước bão to gió lớn nhưng ít ra vẫn còn có sóng bốn G để lướt mạng, còn ở cái xó đảo hoang này thì đến một vạch sóng cũng khum có-nanora!!''

Subaru ôm trán gầm lên trong tiếng gió rít: ``Vấn đề lúc này không phải là sóng điện thoại, mà là làm sao sống sót để rời khỏi cái đảo này kìa!!''

Thế nhưng, vừa dứt lời than khóc, nàng Công chúa Kẹo bỗng nhiên bật dậy như một chiếc lò xo. Cô nàng thoăn thoắt rút chiếc gậy tự sướng gắn điện thoại thông minh ra, quay ngoắt về phía chúng tôi, hớn hở đổi giọng: ``Trước tiên thì ta cứ chụt vài bức ảnh và chỉnh xửa nó đi đã-nora!''

\img{13-11b}

Tôi cạn lời nhìn cô nàng ngây thơ ấy: ``Đến cái hoàn cảnh sinh tử này mà em vẫn còn tâm trạng selfie sống ảo được thì anh cũng đến lạy em rồi đấy, Luna-tan ạ\dots''

Từ cụm loa trên cổ tay tôi, giọng trợ lý ảo Game cất lên đầy gắt gỏng: ``Đúng là bản tính khó dời của bọn con gái sống ảo thời nay. Đứng trước thềm thảm họa diệt vong mà điều đầu tiên nghĩ đến lại là số lượt thả tim trên mạng xã hội.''

Luna nghe thấy liền chu môi hờn dỗi: ``Hai ngừi cằn nhằn quá á-nanora! Thế rốt cục có chụp với tụi này hơm?!''

Tôi giơ tay từ chối dứt khoát: ``Thôi khỏi, anh xin kiếu.''

Naomi đứng bên cạnh tôi khẽ giơ chiếc máy ảnh cơ nhỏ lên, mỉm cười nhẹ nhàng: ``Em chụp cho Luna-oneechan nhé?''

Luna bĩu môi: ``Hơm chụp thì thôi, nàm gì mà căng na\~{} Cảm ơn Naomi-chan nhen! Subarucha-senpai, Suisei-senpai! Hai chị thì sao na? Vào đây làm một kiểu kỷ niệm đi nora!''

Nàng tóc hồng nhiệt tình kéo hai tiền bối của mình lại gần để làm một kiểu ảnh tự sướng. Thế nhưng, trái ngược hoàn toàn với sự phấn khởi của Luna, cả Suisei lẫn Subaru lúc này đều chẳng hề thoải mái chút nào. Sau chuyến đi tàu lắc lư dữ dội giữa tâm bão, khuôn mặt của hai cô nàng trông chẳng khác nào những kẻ mất ngủ kinh niên: mắt quầng thâm sì, mặt mày xanh xao nhợt nhạt, cơ miệng cứng đờ đến mức không thể gượng ép nhếch nổi một nụ cười.

\img{13-11c}

*TÁCH!*

Naomi đứng bên cạnh cũng nhanh tay bấm máy, ghi lại trọn vẹn hậu trường thảm hại của hai tiền bối. Cô bé nhìn vào màn hình máy ảnh rồi khẽ thì thầm với tôi: ``\vie{Anh Kuon ơi\dots\ Trông chị Suisei với chị Subaru giống hệt mấy con gấu trúc bị ốm vậy ạ.}''

Tôi khẽ cười trừ: ``\vie{Cái này người ta gọi là chụp ảnh bất chấp số phận đấy em gái à\dots}''

Bất chợt\dots

*BÍP BÍP BÍP! BÍP BÍP BÍP!!*

Những tiếng còi báo động khẩn cấp rú lên liên hồi từ chiếc đồng hồ đeo tay của tôi, phá tan bầu không khí gượng gạo. Đốm sáng đỏ rực nhấp nháy điên cuồng trên màn hình radar.

Game cất giọng hối thúc, âm điệu vô cùng nghiêm trọng: ``Kuon-user! Báo động đỏ cấp cực hạn!''

Tôi giật mình siết chặt lấy bả vai Naomi: ``Có chuyện gì thế Game!?''

Vị trợ lý ảo đáp: ``Phân tích dữ liệu vệ tinh khí tượng cho thấy: Đây không phải là một cơn bão nhiệt đới thông thường! Vùng dị thường áp thấp này đang sản sinh ra một chuỗi liên hoàn gồm đúng 2022 cơn lốc xoáy mini di chuyển theo quỹ đạo zíc zắc ngoài khơi! Chúng ta phải rời khỏi bờ biển này ngay lập tức!''

Tôi nhìn ra mặt biển đang cuộn trào sóng dữ: ``Rời đi kiểu gì cơ chứ!? Thời tiết này thì làm gì có con tàu nào dám cập bến hòn đảo này nữa!''

Suisei loạng choạng bước tới: ``Chuyện gì đang xảy ra vậy, Kuon-senpai!?''

Tôi và Suisei cùng nhìn vào màn hình quét radar ba chiều của Game. Trên đó, một biểu tượng vòi rồng khổng lồ mang mã số ``Lốc xoáy số 001 / 2022'' đang lao vùn vụt về phía hòn đảo với tốc độ kinh hoàng.

Tôi thốt lên: ``\textbf{Ôi định mệnh!!}''

Trong khi đó, nàng Công chúa Kẹo mắt vẫn dán chặt vào màn hình điện thoại, thoăn thoắt lướt ngón tay chỉnh sửa ảnh mà chẳng hề mảy may để ý tới sự đời: ``Giờ thì zùng filter nào cho da dẻ hồng hào lên đây nora\~{}''

Subaru thảng thốt gào lên: ``\textbf{Luna!! Chạy mau!!}''

Luna vẫn cắm cúi gõ phím, càu nhàu vọng lại: ``Từ từ đã, đang bựn-nora!''

Tôi không thể chịu nổi nữa, dồn hết sức bình sinh hét toáng lên: ``\textbf{Luna! Nhìn lên trời mau!!}''

Tiếng hét như sấm truyền của tôi khiến cô nàng giật bắn mình buông rơi chiếc gậy tự sướng xuống cát. Và khoảnh khắc ngước mắt nhìn lên, toàn bộ tròng mắt của Luna mở to đến cực hạn trong nỗi khiếp đảm tột cùng.

\img{13-11d}

Subaru run rẩy chỉ tay: ``L-Lốc xoáy!! Một cơn lốc xoáy khổng lồ đang lao thẳng vào chúng ta kìa!!''

Cột lốc xoáy đen ngòm cao hàng trăm mét rít lên từng hồi kinh thiên động địa, càn quét mặt biển và lao thẳng vào bờ cát với tốc độ xé gió.

Tôi vội vàng ôm chặt Naomi vào lòng, gập người che chắn: ``Tất cả nằm rạp xuống-- Á! Sui-chan!!''

Suisei hét lên thất thanh: ``\textbf{Aaahhhh!!}''

Game rú lên qua loa: ``Chết tiệt! Cơn lốc số 001 đã cuốn Hoshimachi-san lên trời rồi!!''

Chưa kịp để bất kỳ ai kịp phản ứng hay tìm chỗ bám víu, luồng cuồng phong khủng khiếp đã ập tới như một chiếc máy hút bụi khổng lồ.

*ÙNG ÙNG ÙNG!!*

Như một lẽ tất yếu, toàn bộ tôi, Naomi, Subaru và Luna đều bị lực hút cực mạnh của vòi rồng nhấc bổng lên không trung!

Chúng tôi cứ thế xoay mòng mòng giữa tâm bão, bị những luồng xoáy không khí đưa lên tận tầng mây đen kịt rồi cuốn phăng đi về một phương trời vô định nào đó ngoài đại dương xa xôi.

Giữa tiếng gió rít gào xé rách màng nhĩ và những giọt mưa quất rát rạt vào mặt, tiếng la hét thất thanh của tất cả mọi người cứ thế vang vọng khắp vòm trời xám xịt\dots

\ct

\begin{center}
  \uline{???}
\end{center}

``Ugh\dots\ Chuyện quái gì\dots\ vừa xảy ra thế này\dots''

Tôi khẽ dụi mắt, rồi từ từ mở hàng mi nặng trĩu ra.

Đập vào mắt tôi không còn là mây đen hay giông bão mịt mù nữa, mà lại là một vòm trời xanh ngắt với lác đác vài cụm mây trắng lững lờ trôi, cùng ánh nắng chói chang rọi thẳng vào mặt.

Game cất tiếng thở dài trên cổ tay tôi, giọng điệu vẫn còn hơi chập chờn sau cú va đập: ``Tỉnh rồi à, Kuon-user\dots\ Thật kỳ diệu là các cậu vẫn chưa bị lực ly tâm bẻ gãy cái xương nào.''

Tôi giật mình quay sang bên cạnh. Naomi đang nằm cuộn tròn trên cát, hai tay vẫn ôm chặt lấy chiếc máy ảnh cơ trước ngực. Tôi vội đỡ em gái ngồi dậy: ``Naomi, em có sao không!? Có bị thương ở đâu không!?''

Cô bé chớp chớp mắt, lắc nhẹ đầu rồi kiểm tra thân máy ảnh: ``Em không sao ạ\dots\ Máy ảnh vẫn còn nguyên vẹn, ống kính không bị vỡ, Onii-chan\dots''

Tôi thở phào nhẹ nhõm. Nhìn sang bên cạnh, Subaru cũng vừa mới lồm cồm bò dậy, hai tay ôm đầu lắc lắc cho cát rơi ra khỏi tóc: ``U u\dots\ Ta đang ở cái chốn nào thế này\dots?''

\img{13-11e}

Trước khi tôi kịp mở miệng đáp lời, một giọng nói quen thuộc mang theo tiếng bước chân lạo xạo trên cát đã vang lên ngay bên cạnh: ``Dường như là chúng ta đã bị thổi bay tới một hòn đảo hoang khác rồi.''

Tôi ngước lên nhìn nàng Sao chổi đang phủi cát trên váy: ``À, may quá Sui-chan, em cũng hạ cánh an toàn ở đây rồi\dots\ Ôi cái đầu tôi, nhức buốt cả thái dương\dots''

Tôi nắm lấy tay nàng Vịt, cùng cô nàng kéo nhau đứng dậy.

Đúng lúc đó, một giọng nói trẻ con ngọng nghịu quen thuộc bất ngờ phát ra từ ngay dưới chân chúng tôi: ``Em có thể chụt vài bức ảnh ở đây đực nè-nora☆ !''

Cả bốn người chúng tôi đồng loạt cúi đầu nhìn xuống.

Ở ngay giữa đống cát lún, Luna đang từ từ trồi cái đầu tóc hồng lên như một củ cải non vừa mới mọc. Trên tay cô nàng vẫn nắm chặt chiếc điện thoại thông minh, vẻ mặt tươi rói hớn hở giơ lên tự sướng.

Game cất giọng than thở đầy mỉa mai: ``Tôi thực sự bái phục và ngạc nhiên trước cái thái độ dửng dưng coi trời bằng vung của cô nàng này đấy. Rơi tự do từ độ cao ba nghìn mét xuống mà việc đầu tiên làm vẫn là selfie.''

Tôi, Suisei và Subaru đồng loạt gật đầu tán thành không chút do dự. Naomi thấy vậy liền hướng ống kính máy ảnh về phía Luna, bấm *TÁCH* một cái chụp lại khoảnh khắc trồi lên từ cát đầy hài hước ấy.

\img{13-11f}

Bất chợt, khi đang mải miết lướt ngón tay kiểm tra lại bức ảnh vừa chụp, nàng Công chúa Kẹo bỗng chớp chớp mắt, nhìn thấy một chi tiết kỳ lạ nằm ở góc bãi cát phía sau lưng: ``Ơ? Ai vít cái chữ ét-ô-ét kia zậy nora? Sui-chan hả?''

Suisei lắc đầu: ``Không phải chị. Chị vừa mới tỉnh dậy chưa đầy hai phút thôi.''

Luna ngước mắt nhìn tôi: ``Kuoncha-senpai?''

Tôi xua tay: ``Anh còn đang bận xem Naomi có bị trầy xước gì không, rảnh đâu mà vẽ. Chắc là Game làm bằng cảm biến nhiệt chăng?''

Game phủ nhận ngay tắp lự: ``Đừng vu oan cho tôi. Tôi không có tay chân để đi cào cát viết chữ. Oozora-san thì lại càng không thể rồi.''

Cả đám chúng tôi đứng trầm ngâm suy nghĩ xem dòng chữ cầu cứu bí ẩn kia từ đâu ra, hoàn toàn quên bẵng đi hoàn cảnh éo le của bản thân hiện tại.

Phải đến khi Subaru giật nảy mình hét toáng lên, mọi người mới sực tỉnh khỏi cơn mê: ``Nguy to rồi!! Hòn đảo này không một bóng người!! Ta phải đi tìm kiếm sự giúp đỡ thôi!!''

Suisei nhíu mày: ``Nhưng mà tìm bằng cách nào cơ chứ?''

Tôi nhìn quanh bờ biển hoang vu: ``Đến việc đây là đảo nào trên bản đồ chúng ta còn chẳng biết nữa là\dots''

Game cất giọng vô cảm: ``Cơ sở dữ liệu bản đồ ngoại tuyến của tôi chưa cập nhật tọa độ vùng biển hẻo lánh này, nên là tôi chịu chết.''

Đúng lúc đó, cô nàng mê chụp ảnh tự sướng bỗng nhiên vung gậy selfie lên, hăng hái xung phong: ``Dễ ợt mà-nora! Chúng ta cứ chia za rồi đi tìm quanh đảo là đực-nora!''

Dứt lời, chẳng thèm đợi ai phê duyệt hay bàn bạc kế hoạch, cả ba cô nàng Suisei, Subaru và Luna đồng loạt cắm đầu cắm cổ chạy thục mạng về ba hướng khác nhau, lao thẳng vào trong bụi cây rậm rạp giữa đảo.

\img{13-11g}

Tôi đứng chôn chân tại chỗ, hốt hoảng hét với theo: ``\textbf{Ấy khoan đã! Đừng có tản ra như thế!! Nguy hiểm lắm đó các em ơi!!}''

Thế nhưng, đáp lại lời cảnh báo khản cổ của tôi chỉ là tiếng lá cây xào xạc. Cả ba bóng người đã hoàn toàn biến mất sau những rặng dừa hoang dại.

Và rồi, như một định luật bất thành văn của cuộc đời, ngay giữa bầu trời phía trên hòn đảo bỗng nhiên hiện ra một dòng chữ to đùng in đậm:

\begin{center}
  \textbf{【TẤT CẢ ĐỀU BỊ LẠC】}
\end{center}

Tôi ôm trán lẩm bẩm, mặt méo xệch: ``Cái bảng chữ thông báo quái thai từ trên trời rơi xuống này là sao chứ? Rốt cuộc cái hòn đảo này bị cái gì ám vậy trời\dots?''

Naomi kéo nhẹ vạt áo tôi: ``Onii-chan\dots\ Đồng hồ của anh lại kêu nữa kìa.''

*BÍP BÍP BÍP BÍP!!*

Âm thanh báo động tử thần lại một lần nữa réo vang rợn người.

Game thảng thốt phát loa: ``Cảnh báo khẩn cấp! Radar phát hiện luồng xoáy đối lưu cực mạnh! Cơn lốc xoáy số 412 trong chuỗi 2022 cơn lốc đang di chuyển thẳng tới hòn đảo này với vận tốc hai trăm kilomet trên giờ!!''

Tôi còn chưa kịp chửi thề lấy một câu thì từ phía chân trời, một cơn lốc xoáy khổng lồ khác đã cuốn phăng rặng cây, quét sạch mọi bụi rậm trên đảo, cuốn phăng cả ba cô nàng vừa chạy vào rừng lẫn tôi và Naomi lên thẳng chín tầng mây.

Những tiếng gào thét kinh hoàng lại một lần nữa xé toạc bầu trời nhiệt đới giữa ban ngày ban mặt\dots

\ct

\begin{center}
  \uline{???}
\end{center}

Cũng giống hệt như kịch bản lần trước, tôi và Subaru lại lồm cồm chống tay bò dậy khỏi bờ cát.

Thứ đầu tiên đập vào mắt chúng tôi lúc này là một vòm trời đã chuyển dần sang sắc cam rực rỡ của buổi chiều tà. Hoàng hôn buông xuống trên mặt biển nhuốm một màu đỏ quạch.

Game cất tiếng hỏi thăm đầy uể oải: ``Hai người lại tỉnh rồi đấy à?''

Tôi thở hắt ra một hơi, phủi cát dính trên má Naomi: ``Ừ\dots\ Vẫn còn nguyên vẹn mạng sống. Cơ mà\dots''

Subaru ôm lấy hai thái dương, giọng nói run rẩy: ``Lại là cái xó xỉnh quỷ quái nào nữa đây trời\dots!?''

Và đúng như một đoạn băng ghi hình được tua lại, Suisei lại từ đằng xa lững thững bước tới giải đáp thắc mắc của chúng tôi bằng chất giọng bình thản đến phát sợ: ``Lại một hòn đảo hoang khác nữa rồi, chắc vậy.''

\img{13-11h}

Tôi ngửa mặt lên trời than thở: ``Trời đất quỷ thần ơi\dots\ Sao hôm nay lại có nhiều lốc xoáy đến mức phi lý thế này cơ chứ\dots''

Bất chợt, ở một góc bãi cát cách đó không xa, Luna - người vừa mới đưa điện thoại lên định đăng tấm ảnh tự sướng chụp hoàng hôn vừa xong - bỗng nhiên thét lên một tiếng kinh hoàng đến xé lòng.

Tôi giật bắn mình, vội vàng kéo Naomi chạy tới bên cạnh cô nàng: ``C-Chuyện gì thế Luna-tan!? Lại có quái vật biển hay lốc xoáy nữa à!?''

Nàng Công chúa xứ Kẹo ngước khuôn mặt nhạt nhòa nước mắt lên nhìn tôi, chỉ vào chiếc điện thoại trên tay: ``Em khum thể tải ảnh lên mạng đực-nora!!''

Chúng tôi nhìn vào màn hình. Bức ảnh selfie của Luna với biểu cảm chu môi đáng yêu vẫn đang bị một chiếc vòng tròn xoay xoay (loading spinner) quay mòng mòng vô tận giữa màn hình.

\img{13-11i}

Subaru sực nhận ra vấn đề cốt lõi, thốt lên: ``Không lẽ\dots\ hòn đảo này hoàn toàn nằm ngoài vùng phủ sóng viễn thông rồi sao!?''

Game cất giọng ngán ngẩm: ``Chả trách nãy giờ định vị vệ tinh của tôi hoàn toàn mất tín hiệu. Chúc mừng các cậu, chúng ta đang ở một nơi mà sóng điện thoại còn không thèm ghé qua.''

Suisei khoanh tay thở dài: ``Đó thực sự là một vấn đề lớn đấy\dots''

Và thế là, thảm kịch thực sự bùng nổ.

Hiểu ra được rằng bức ảnh sống ảo tâm huyết của mình không thể khoe lên mạng xã hội, nàng Công chúa Kẹo lập tức lăn đùng ngã ngửa ra bãi cát, lăn lộn mấy vòng ăn vạ: ``\textbf{Nnnaaaaa!! Hơm chịu đâu!!}''

Rồi cô nàng bật dậy, quỳ gối trên cát, hai tay đập bồm bộp xuống nền đất trong nỗi uất ức tột cùng: ``Em sẽ zít thời zan chờ đợi như thế nào nữa đây-nora\~{}!?''

\img{13-11j}

Subaru, tôi và Game đồng loạt đưa mắt nhìn cô nàng với vẻ mặt không thể nào cạn lời hơn.

Game lạnh lùng châm chọc: ``Ở giữa bờ cõi hoang vu sống chết chưa rõ mà cô vẫn còn dư thừa nơ-ron thần kinh để nghĩ ra cách giết thời gian cơ đấy\dots''

Tôi khoanh tay lắc đầu bất lực: ``Đến nước này thì anh cũng xin giơ hai tay đầu hàng em rồi Luna ạ. Sống ảo vừa vừa phải phải thôi chứ.''

Naomi đứng bên cạnh giơ máy ảnh bấm *TÁCH* thêm một kiểu, chụp lại cận cảnh khuôn mặt mếu máo ăn vạ của nàng Công chúa, rồi nhỏ nhẹ bảo: ``Luna-oneechan đừng khóc nữa mà\dots\ Em chụp cho chị bằng máy ảnh phim này, đẹp lắm đó ạ.''

Luna dẩu môi phụng phịu: ``Hơm chịu đâu\~{}! Em mún được các bạn fan trên Twitter bít em đang nàm gì-noraa!''

Subaru nổi đóa quát lớn: ``\textbf{Đừng có nằm lăn ra ăn vạ như con nít nữa! Không có mạng thì ta đi tìm chỗ có sóng! Dù gì đi nữa, ta nên tách nhau ra để tìm cột thu phát tín hiệu nào!!}''

Và rồi, lịch sử đen tối lại lặp lại một cách hoàn hảo: không ai bảo ai, cả ba cô nàng thầ tượng lại tiếp tục chia làm ba hướng, cắm đầu cắm cổ lao vùn vụt vào rừng rậm nhiệt đới để tìm sóng điện thoại, tiếp tục bỏ mặc tôi, Naomi và Game ở lại bờ biển.

\img{13-11k}

Tôi ôm trán thở dài thườn thượt: ``Đã bảo đó không phải là cách giải quyết khôn ngoan rồi mà\dots''

Và dĩ nhiên, dòng chữ thông báo chết tiệt quen thuộc lại hiện ra lơ lửng giữa trời:

\begin{center}
  \textbf{【TẤT CẢ ĐỀU BỊ LẠC】}
\end{center}

Ngay giây tiếp theo\dots

*BÍP BÍP BÍP BÍP BÍP!!*

Game rú lên điên cuồng qua cụm loa: ``Cảnh báo va chạm khẩn cấp! Cơn lốc xoáy số 1314 trong chuỗi 2022 cơn lốc đang ập tới ngay sau lưng các cậu!!''

Tôi chỉ kịp thốt lên một câu chửi thề ngắn gọn: ``\textbf{Mẹ kiếp cái thời tiết năm 2022 này!!}''

*ÙNG ÙNG ÙNG!!*

Cơn lốc xoáy hoàng hôn màu cam đỏ khổng lồ lại quét sạch bãi cát, nuốt chửng cả năm con người chúng tôi lên không trung lần thứ ba trong ngày.

Giữa những tiếng thét gào khản đặc cả cổ họng, chúng tôi lại bị lực xoáy cuốn bay về một phương trời xa xôi nào đó\dots

\ct

\begin{center}
  \uline{???}
\end{center}

Một lần nữa, tôi, Naomi và Subaru lại chồm dậy khỏi mặt đất.

Trải qua ba lần bay lượn trên không trung như chơi tàu lượn siêu tốc, dường như tiền đình và cơ thể chúng tôi đã dần hình thành một cơ chế miễn dịch kỳ lạ với việc bị lốc xoáy cuốn bay.

Bầu trời lúc này đã chuyển sang màu xanh thẫm của buổi chập tối. Ánh đèn đường le lói từ xa hắt lại khiến tôi nhận ra một điều kỳ lạ: mặt đất dưới chân không còn là cát biển nữa, mà là một bờ kè bê tông kiên cố.

Ở phía trước, Luna đang giơ cao chiếc điện thoại lên trời, hai mắt dán chặt vào cột sóng mạng.

\img{13-11l}

Tôi phủi bụi trên áo cho Naomi, chép miệng đoán trước: ``Để anh đoán xem nhé Subaru, em sẽ lại ngơ ngác hỏi\dots''

Subaru chớp mắt, buột miệng theo đúng phản xạ: ``Ể!? Ta lại đang ở cái xó nào đây\dots?''

Game cất giọng đều đều: ``Và rồi Hoshimachi-san sẽ lại lững thững bước ra trả lời\dots''

Suisei từ sau bức tường bê tông bước ra, giơ tay chào tỉnh bơ: ``Có vẻ như lần này là bờ biển phía nam vùng Kanto rồi đấy.''

Vị Trợ lý ảo lập tức khựng lại: ``Khoan đã, cái gì cơ!? Lần này về được đất liền thật rồi á!?''

Câu trả lời bất ngờ của nàng Sao chổi khiến tất cả chúng tôi trố mắt kinh ngạc: ``Thật á Sui-chan!?''

Đúng lúc đó, nàng Công chúa tóc hồng bỗng nhiên nhảy cẫng lên reo hò trong sung sướng tột cùng, hai mắt sáng rực lên như hai vì sao tinh tú: ``Aaaa!! Em bắt đực tín hịu bốn G rùi nè-nora!!''

\img{13-11m}

Chúng tôi lập tức nhào tới, dán mắt vào ứng dụng bản đồ số mà Luna vừa mở ra trên màn hình.

Trên chấm xanh định vị GPS, dòng chữ địa danh hiển thị rõ mồn một: Bờ biển Chigasaki, khu vực Shounan, tỉnh Kanagawa, thuộc vùng Kanto.

Tất cả chúng tôi đồng loạt thở phào một hơi nhẹ nhõm đến run cả người: ``Chả trách sao lại có đường nhựa và đèn đường\dots\ Về được tới đất liền rồi!!''

Ít nhất thì\dots\ cơn nghiện sống ảo và chiếc điện thoại thông minh của Luna cuối cùng cũng phát huy được một công dụng có ích cho nhân loại.

%==========================================TRUYỆN PHỤ=========================================
\omk

Subaru ngồi bệt xuống bờ kè bê tông, xoa xoa bắp chân mỏi nhừ: ``Giờ ta đã biết được vị trí chính xác rồi\dots''

Suisei nhìn quanh đại lộ vắng ngắt không một bóng xe cộ: ``Nhưng mà bão gió thế này, toàn bộ tàu điện ngầm lẫn xe buýt liên tỉnh đều đã ngừng hoạt động hết rồi. Ta về văn phòng bằng cách nào đây?''

Trong khi hai bậc tiền bối đang vò đầu bứt tai lo lắng tìm phương tiện di chuyển, thì ở phía mép nước, Luna vẫn điềm nhiên tạo đủ mọi dáng uốn éo, giơ hai ngón tay chữ V tự sướng để kịp đăng bức ảnh hoàng hôn vừa tải xong lên trang cá nhân, hoàn toàn ngó lơ thực tế nghiệt ngã.

Naomi đứng bên cạnh tôi, khẽ giơ chiếc máy ảnh lên ngắm về phía rặng mây đen đang ùn ùn kéo tới từ đường chân trời: ``Onii-chan\dots\ Gió lại bắt đầu rít mạnh hơn rồi kìa ạ.''

Đúng như lời cô bé nói, những tiếng còi báo động chát chúa quen thuộc lại một lần nữa réo vang inh ỏi trên chiếc đồng hồ của tôi.

Tôi day day trán, buông một tiếng thở dài trong sự ngán ngẩm tột cùng: ``Lại có lốc xoáy nữa đúng không Game?''

Từ cổ tay tôi, Game thở hắt ra một hơi, giọng điệu như thể đã buông xuôi chấp nhận số phận: ``Chính xác là như vậy. Cơn lốc xoáy mang mã số 2022 - cơn lốc cuối cùng khép lại chuỗi siêu dị tượng khí tượng năm nay - đang tiến thẳng vào bờ biển Chigasaki với tốc độ ba trăm kilomet trên giờ.''

Dứt lời, chẳng kịp để chúng tôi kêu lên một tiếng than khóc nào, một cơn lốc xoáy đen ngòm khổng lồ mang sức mạnh kinh thiên động địa đã quét sạch toàn bộ bờ biển Chigasaki.

*ÙNG ÙNG ÙNG!!*

Cây cối ven đường, biển báo giao thông và cả năm con người chúng tôi lại một lần nữa bị hút bổng lên tầng bình lưu, xoáy tít mù trong cơn cuồng phong dữ dội nhất của đêm hè.

Giữa bầu trời đêm mù mịt đầy giông bão, cơn lốc xoáy số 2022 càn quét thẳng từ bờ biển Kanagawa bay vút về hướng vịnh Tokyo\dots

\ct

\begin{center}
  \uline{???}
\end{center}

Tại Văn phòng chính của \hll.

Một lần nữa, tôi và Subaru lại lờ đờ chớp mắt tỉnh dậy.

Thế nhưng lần này, chúng tôi lập tức cảm nhận được một sự khác biệt rõ rệt ở lưng: cảm giác êm ái, mịn màng của tấm thảm trải sàn văn phòng quen thuộc, chứ không còn là sỏi đá hay cát biển bụi bặm như những lần trước nữa.

Subaru xoa xoa đầu, ngơ ngác nhìn quanh: ``\dots Ta lại đang ở đâu đây nữa thế này?''

Suisei đứng khoanh tay bên cạnh chiếc bàn làm việc, nhếch môi cười: ``Về tới văn phòng công ty rồi.''

Tôi và Subaru giật bắn mình bật dậy. Đúng là khung cảnh văn phòng \hll\ quen thuộc ở Thành phố Điện tử. Cơn lốc xoáy số 2022 bằng một cách thần kỳ khó tin nào đó đã đóng vai trò như một dịch vụ giao hàng siêu tốc, ném thẳng chúng tôi xuyên qua cửa sổ tầng cao rơi an toàn xuống phòng khách!

Tôi vội quay sang kiểm tra Naomi: ``Naomi! Em có sao không!?''

Cô em út ngồi dậy trên chiếc ghế sô-pha mềm mại, ôm máy ảnh mỉm cười rạng rỡ: ``Em không sao đâu Onii-chan! Chúng ta về đến nhà thật rồi nè!''

Thế nhưng, niềm vui trở về còn chưa kịp nở trọn thì một luồng gió bão kinh hoàng từ bên ngoài bất ngờ thốc mạnh qua khung cửa sổ vỡ nát, làm cả đám chúng tôi loạng choạng suýt ngã nhào xuống sàn.

Game cất giọng ngán ngẩm qua cụm loa: ``\dots Tôi quên béng mất một điều hiển nhiên: Tâm bão năm 2022 hôm nay đổ bộ trực diện vào Thành phố Điện tử.''

Suisei gật đầu: ``Ừ ha\dots''

Subaru vội vã hô hoán: ``Mau dậy đóng cửa sổ lại đi mấy người ơi!! Gió giật tung hết đồ đạc bây giờ!!''

Nàng Ca sĩ chỉ tay xuống sàn nhà, dở khóc dở cười: ``Chị cũng muốn đóng lắm chứ Subaru, nhưng mà em nhìn xem\dots''

Tôi nhìn theo ngón tay cô nàng. Dưới sàn nhà, toàn bộ khung kính cửa sổ đã bị cơn lốc xoáy thổi vỡ nát tan tành từ lúc nãy. Hàng trăm tờ tài liệu, giấy tờ hợp đồng bay lượn tung tóe khắp phòng như bướm lượn, ngoài trời thì mưa xối xả như trút nước bạc cả một khoảng không.

Subaru nhìn quanh quất, bỗng giật mình thảng thốt: ``Ủa khoan đã\dots\ Luna đâu rồi!?''

Cả căn phòng lập tức chết lặng. Không thấy bóng dáng nàng Công chúa Kẹo đâu cả! Chẳng lẽ cô nàng bị rơi rụng lại ở bờ biển Kanagawa rồi sao!?

Chưa kịp để Subaru tá hỏa tam tinh gọi điện báo cảnh sát, một tiếng chuông thông báo *ting* vang lên giòn giã từ chiếc điện thoại của cô nàng Vịt.

Màn hình sáng lên thông báo có một dòng tweet mới nhất từ tài khoản Twitter của Luna.

Subaru cuống cuồng mở ra xem. Ngay giây tiếp theo, cả tôi, Subaru, Suisei, Naomi và Game đều đồng loạt trợn trừng hai mắt, há hốc mồm kinh ngạc đến mức chết trân tại chỗ!

Đập vào mắt chúng tôi là một bức ảnh selfie tự sướng với độ phân giải cực nét: Luna đang nở một nụ cười tươi rói, giơ hai ngón tay làm dấu chữ V, hai bím tóc hồng bay phấp phới giữa không trung khi cô nàng đang\dots\ bay lượn lờ ở độ cao hàng ngàn mét bên trong chiếc phễu của cơn lốc xoáy!

Kèm theo đó là dòng trạng thái hồn nhiên đến rợn tóc gáy:

\txtbx{``Bay cao chên này zui wá-nanora\~{}!''}

Subaru ôm chặt lấy đầu, gào thét trong tuyệt vọng: ``\textbf{L-Luna!?!??}''

Suisei ôm trán lắc đầu cạn lời: ``Thế quái nào mà cơn lốc ném bốn người chúng ta rơi chuẩn xác vào văn phòng, mà em ấy vẫn còn có thể kẹt lại trên trời để chụp ảnh sống ảo được nhỉ!?''

Tôi thở dài thườn thượt, nhìn cô em út Naomi đang bật cười khúc khích bên cạnh: ``Đúng là sống ảo bất chấp cả sinh mạng mà\dots''

Game lạnh lùng buông lời tổng kết: ``Hệ thống phân tích tâm thần học của tôi xin chính thức từ chối tiếp nhận trường hợp của Himemori-san. Đây là một ca bệnh nan y vô phương cứu chữa.''

Đến cuối cùng, sau khi cơn lốc xoáy số 2022 tan dần trên bầu trời Tokyo, bằng một phép màu kỳ diệu nào đó của thế giới cổ tích, nàng Công chúa Kẹo cũng đã đáp cánh an toàn xuống ban công văn phòng mà không hề xây xước lấy một sợi tóc. Dĩ nhiên, vừa đặt chân xuống đất, cô nàng đã phải hứng chịu một trận lôi đình mắng mỏ té tát suốt nửa tiếng đồng hồ từ Subaru, khiến tôi và Naomi phải đứng ra can ngăn mãi mới êm thấm.

Tôi nhìn đống giấy tờ bừa bộn khắp sàn và khung cửa sổ toang hoác đang đón từng đợt gió bão, thầm lắc đầu cảm thán trong lòng: chưa bao giờ có một chuyến đi nào lại đầy bão táp theo đúng cả nghĩa đen lẫn nghĩa bóng như thế này, kể từ sau chuyến đi biển đầy sóng gió của gia đình tôi hồi tháng trước!

\begin{center}
  {\ab BONUS GIF}
  \gif{13-11n}{1}{43}
\end{center}

\end{document}